Our scheme uses the MPC to generate optimal control inputs while the Neural part is used to model the dynamics and kinematics of our system.
The Unicycle has been modeled with its kinematic model while for the quadrotor we have used its dynamic model.
\subsection{Unicycle kinematic model}
The kinematic model of the unicyle is a simplified representation of the movement of the vehicle
in terms of position and orientation. This simple model has been derived directly from the so-called "rolling without slipping constraint". 
The state space is represented by the robot cartesian position and orientation $[x, y, \theta]$, while the input actions are represented by the driving and steering velocities the vehicle can assume $[v, \omega]$.
The kinematic model of a univcyle robot it is described by the following equations:
\begin{align}
    \label{eq:unicycle}
    \begin{cases}
        \dot{x}(t) = v(t) \cos(\theta(t)) \\
        \dot{y}(t) = v(t) \sin(\theta(t)) \\
        \dot{\theta}(t) = \omega(t) \\
    \end{cases}
\end{align} 
\subsection{Quadrotor dynamic model}
For the quadrotor, we have used its dynamic model. We can describe this second order model by starting from the two following equations:
\begin{align}
    \label{eq:uav_start}
    \begin{cases}
         m \ddot{p} = f_{t} + f_{g} + f_{a} + f_{d} \\
         J \dot{\omega} + \omega \times (J \omega) = \tau
    \end{cases}  
\end{align}
where the first equation, which is derived from Newton's Second Law, represents the Translational Dynamics while the second one, derived from Euler's equations, represents the Rotational Dynamics.
In the first equation $f_{t}$ represents the total thrust force: $f_{t} = R \cdot \begin{bmatrix} 0 & 0 & T \end{bmatrix} ^T$, where R is the rotation matrix of $(\phi, \theta, \psi)$, the roll, pitch and yaw angles. The total thrust force T produced by the rotors is defined by the following equation:
\begin{equation}
    \label{eq:thrust}
    T = f_1 + f_2 + f_3 + f_4,\quad \text{with} \quad f_i = k_f \omega_i^2
\end{equation}
where $k_f$ is the thrust coefficient and $\omega_i$ is the angular speed of the $i$-th rotor.
For the Rotational Dynamics, $\tau$ represents the torques in the body frame and can be expressed as follows:
\begin{equation}
    \tau =
    \begin{bmatrix}
        \tau_{\phi} \\
        \tau_{\theta} \\
        \tau_{\psi}
    \end{bmatrix}
    =
    \begin{bmatrix}
        L(f_1 - f_3) \\
        L(f_2 - f_4) \\
        \tau_{\psi}
    \end{bmatrix}
\end{equation}
where $L$ is the distance from the center of mass of each rotor.


Given the previous equations, if we assume aereodynamics and gyroscopic effects to be negligible and that we have small roll and pitch angles, we can derive the simplified model fo the quadrotor:
\begin{align}
    \label{eq:uav_second_order}
    \begin{cases}
        \ddot{x} = -(\cos\psi \sin\theta \cos\phi + \sin\psi \sin\phi) \frac{T}{m}\\
        \ddot{y} = -(\sin\psi \sin\theta \cos\phi - \cos\psi \sin\phi) \frac{T}{m} \\
        \ddot{z} = -(\cos\theta \cos\phi) \frac{T}{m} +g\\
        \ddot{\phi} = \frac{\tau_{\phi}}{I_{x}}\\
        \ddot{\theta} = \frac{\tau_{\theta}}{I_{y}}\\
        \ddot{\psi} = \frac{\tau_{\psi}}{I_{z}}\\ 
    \end{cases}  
\end{align}
